% Part: computability
% Chapter: recursive-functions
% Section: non-pr-functions

\documentclass[../../../include/open-logic-section]{subfiles}

\begin{document}

\olfileid[hi]{cmp}{rec}{npr}
\olsection{गैर-आदिम पुनरावर्ती फलन}

आदिम पुनरावर्ती फलन सहज रूप से संगणनीय समझे जाने वाले सभी फलनों को पूरा
नहीं करते। यह सहज रूप से स्पष्ट होना चाहिए कि हम सभी एक-स्थानी आदिम
पुनरावर्ती फलनों की ऐसी सूची $f_0$, $f_1$, $f_2$,~\dots{} बना सकते हैं
कि $f_x$ का मान निवेश $y$ पर प्रभावी रूप से संगणित किया जा सके; दूसरे
शब्दों में, निम्न प्रकार परिभाषित फलन $g(x,y)$ संगणनीय है:
\[
g(x,y) = f_x(y)
\]
किंतु तब निम्न फलन भी संगणनीय है:
\begin{eqnarray*}
h(x) & = & g(x,x) + 1 \\
& = & f_x(x) +1.
\end{eqnarray*}
प्रत्येक आदिम पुनरावर्ती फलन $f_i$ के लिए $h$ और $f_i$ के मान $i$ पर
भिन्न हैं। अतः $h$ संगणनीय है, किंतु आदिम पुनरावर्ती नहीं; और $g$ के
बारे में भी यही कहा जा सकता है। यह कैंटर के विकर्णीकरण तर्क का
``प्रभावी'' रूप है।

ऐसे संगणनीय फलनों के और अधिक स्पष्ट उदाहरण दिए जा सकते हैं जो आदिम
पुनरावर्ती नहीं हैं। मसलन, संकेत $g^n(x)$ से
$g(g(\dots g(x)))$ निरूपित हो, जिसमें कुल $n$ बार $g$ आता है; और फलनों
का अनुक्रम $g_0,g_1,\dots$ इस प्रकार परिभाषित करें:
\begin{eqnarray*}
g_0(x) & = & x+1 \\
g_{n + 1}(x) & = & g_n^x(x)
\end{eqnarray*}
आप सत्यापित कर सकते हैं कि प्रत्येक फलन $g_n$ आदिम पुनरावर्ती है। हर
अगला फलन अपने पूर्ववर्ती से बहुत अधिक तीव्रता से बढ़ता है; $g_1(x)$,
$2x$ के बराबर है, $g_2(x)$, $2^x \cdot x$ के बराबर है, और $g_3(x)$
लगभग ऐसे घातांकीय स्तंभ की तरह बढ़ता है जिसमें $x$ बार $2$ आता है।
एकरमान--पीटर (Ackermann--P\'eter) फलन मूलतः $G(x) = g_x(x)$ फलन है, और
दिखाया जा सकता है कि यह किसी भी आदिम पुनरावर्ती फलन से अधिक तीव्रता से
बढ़ता है।

आइए आदिम पुनरावर्ती फलनों को परिगणित करने के प्रश्न पर लौटें। याद रखें कि
हमने प्रत्येक आदिम पुनरावर्ती फलन को औपचारिक संकेत दिया है; इसलिए संकेतों
को परिगणित करना पर्याप्त है। हम प्राकृतिक संख्या $\#(F)$ को प्रत्येक
संकेत $F$ के साथ पुनरावर्ती रूप से इस प्रकार जोड़ सकते हैं:
\begin{eqnarray*}
\#(0) & = & \langle 0 \rangle \\
\#(S) & = & \langle 1 \rangle \\
\#(\Proj{n}{i}) & = & \langle 2, n, i \rangle \\
\#(\fn{Comp}_{k,l}[H,G_0,\dots,G_{k-1}]) & = & \langle
3,k,l,\#(H),\#(G_0),\dots,\#(G_{k-1}) \rangle \\
\#(\fn{Rec}_l[G,H]) & = & \langle 4, l, \#(G), \#(H) \rangle
\end{eqnarray*}
यहाँ हम इस तथ्य का उपयोग कर रहे हैं कि पिछले अनुभाग के कूटों से संख्याओं
के प्रत्येक अनुक्रम को प्राकृतिक संख्या माना जा सकता है। निष्कर्ष यह है
कि प्रत्येक कूट को एक प्राकृतिक संख्या दी जाती है। निस्संदेह, कुछ
अनुक्रम (और इसलिए कुछ संख्याएँ) संकेतों के संगत नहीं होते; किंतु हम $f_i$
को उस एक-स्थानी आदिम पुनरावर्ती फलन के रूप में ले सकते हैं जिसका संकेत
$i$ से कूटबद्ध है, यदि $i$ ऐसा कोई संकेत कूटबद्ध करता है; और अन्यथा उसे
अचर $0$ फलन मान सकते हैं। अंतिम परिणाम यह
है कि हमारे पास एक-स्थानी आदिम पुनरावर्ती फलनों को परिगणित करने का स्पष्ट
ढंग है।

(वास्तव में, अचर शून्य फलन जैसे कुछ फलन सूची में एक से अधिक बार आएँगे।
यह केवल हमारे कूटन की कृत्रिम उपज नहीं, बल्कि इस तथ्य का परिणाम भी है कि
अचर शून्य फलन के एक से अधिक संकेत हैं। आगे हम देखेंगे कि इन
पुनरावृत्तियों को संगणनीय रूप से टाला नहीं जा सकता; मसलन, ऐसा कोई
संगणनीय फलन नहीं है जो यह निर्णय करे कि दिया हुआ संकेत अचर शून्य फलन को
निरूपित करता है या नहीं।)

अब हम फलन $g(x,y)$ को $f_x(y)$ से दिया हुआ मान सकते हैं, जहाँ $f_x$
अभी वर्णित परिगणन का संदर्भ देता है। हम कैसे जानते हैं कि $g(x,y)$
संगणनीय है? सहज रूप से यह स्पष्ट है: $g(x,y)$ संगणित करने के लिए पहले
$x$ को ``खोलें'' और देखें कि वह किसी एक-स्थानी फलन का संकेत है या नहीं।
यदि है, तो निवेश~$y$ पर उस फलन का मान संगणित करें।

\begin{tagblock}{TMs}
\begin{digress}
आपको शायद पहले से विश्वास हो गया हो कि (कुछ परिश्रम से!) ऐसा करने वाला
प्रोग्राम (मान लें Java या C++ में) लिखा जा सकता है; और अब हम
चर्च--ट्यूरिंग परिकल्पना का सहारा ले सकते हैं, जो कहती है कि सहज रूप से
संगणनीय हर चीज़ को ट्यूरिंग मशीन से संगणित किया जा सकता है।

निस्संदेह, यह दिखाने का अधिक सीधा ढंग कि $g(x,y)$ संगणनीय है, उसे
संगणित करने वाली ट्यूरिंग मशीन का स्पष्ट वर्णन करना है। विशेषतः इससे
चर्च--ट्यूरिंग परिकल्पना और सहज-बोध के सहारे से बचा जा सकेगा। शीघ्र ही
हमारे पास यह दिखाने के लिए पर्याप्त साधन होंगे कि $g(x,y)$ संगणनीय है;
हम संगणन के ऐसे प्रतिरूप का सहारा लेंगे जिसका ट्यूरिंग मशीन पर
\emph{अनुकरण} किया जा सकता है: अर्थात् पुनरावर्ती फलन।
\end{digress}
\end{tagblock}

\end{document}
